\begin{frame}{ML1 개념지도}
  \begin{columns}[T]
    \begin{column}{0.52\textwidth}
      \centering
      \includegraphics[width=\linewidth]{ch16_3_concept_map.png}
      \vskip0.4em
      {\footnotesize Ch02~15 개념지도 — 4개 박스 + 색 두 줄}
    \end{column}
    \begin{column}{0.46\textwidth}
      \small
      \begin{itemize}
        \item \textbf{4개 박스}: \kb{① 데이터·문제 정의}(행의
          단위/정답 유무/형태) $\to$ \kb{② 모델}(10가지 도구)
          $\to$ \kb{③ 학습}(\textbf{4단계 구조}) $\to$ \kb{④
          평가}(3분할·train fit / val에서만 / test 딱 한 번)
        \item \textbf{붉은 점선}: ④에서 ①로 \textbf{돌아감}
          (과적합·누수면 다시) — ``되돌아갈 수 없는 것은
          \textbf{test 한 번}뿐''
        \item \textbf{금색 점선}: ③에서 \textbf{ML2로 나감} —
          정답 $y$가 \textbf{보상 $r$}로, 데이터가 \textbf{
          에이전트의 경험}으로
      \end{itemize}
    \end{column}
  \end{columns}
  \vskip0.4em
  \begin{block}{쓰기는 단순: Ch02~15의 모델을 하나 골라 ①②③④에
    놓아본다}
    \textbf{GMM}: ③에서 ``\textbf{EM}이 관측 데이터의 \textbf{
    우도}를 최대화하는 방향으로 $\theta$ 조정''.
    \textbf{VAE}: ②에서 ``\textbf{ELBO}가 우도의 \textbf{
    하한}''(exact MLE는 어려우므로).
    \textbf{GAN}: ③에서 ``\textbf{min-max — 게임}'' — 세 개를
    \textbf{전부} 놓을 수 있으면, 이 학기는 정말 끝난 것이다.
  \end{block}
\end{frame}
